% ECONOMICS_PROBLEM: {"id": "D91-609", "title": "Motivated beliefs in economic decisions", "jel_code": "D91", "source_id": "OP-0305"}
% !TeX program = xelatex
\documentclass[11pt,a4paper]{article}
\usepackage[margin=23mm]{geometry}
\usepackage{fontspec}
\setmainfont{Latin Modern Roman}
\usepackage{amsmath,amssymb,mathtools}
\usepackage{xcolor,enumitem,booktabs,longtable,array,xurl,hyperref}
\definecolor{muted}{HTML}{53636C}
\hypersetup{colorlinks=true,urlcolor=blue,linkcolor=blue}
\setlength{\parindent}{0pt}
\setlength{\parskip}{5pt}
\newcommand{\R}{\mathbb R}
\newcommand{\E}{\mathbb E}
\newcommand{\N}{\mathbb N}
\newcommand{\ind}{\mathbf 1}
\newcommand{\Var}{\operatorname{Var}}
\newcommand{\Cov}{\operatorname{Cov}}
\newcommand{\supp}{\operatorname{supp}}
\DeclareMathOperator*{\argmin}{arg\,min}
\DeclareMathOperator*{\argmax}{arg\,max}
\newcommand{\entrymeta}[3]{{\sffamily\small\color{muted}\textbf{\texttt{#1}}\quad|\quad #2\par #3\par}}
\newcommand{\Problemref}[1]{\texttt{#1}}
\newcommand{\JELref}[2]{\texttt{#2} (JEL \texttt{#1})}
\begin{document}
\section*{Motivated beliefs in economic decisions}
\textbf{Problem ID:} \texttt{D91-609}\quad \textbf{JEL:} \texttt{D91}\par
% BEGIN ECONOMICS STATEMENT
\label{problem:OP-0305}
{\sffamily\small\color{muted}Original subject group: Behavioral and experimental economics\par}
\label{definitions:problem:30007587}
\textbf{Audit classification:} Background; proposed extension. Publisher indexing and the authors’ research page verify the 2025 stock-return experiment. Full publisher text remained inaccessible; the general self-image, persistence and savings agenda is editorial.
\subsection*{Definitions and precise research target}
\textbf{Model and research target.} Fix an economic decision with uncertain monetary return \(R\) and objectively specified public information \(I\). Let \(Z_i\in\{0,1\}\) randomly assign an ownership or self-image treatment intended to change the desirability of a favorable return without changing available information. At horizon \(t\), observe an incentivized subjective return forecast \(B_i(Z_i,t)\) and an economic action \(A_i(Z_i,t)\), such as a portfolio share. The belief and action effects in target population \(P\) are
\[\tau_B(t)=\mathbb E_P[B_i(1,t)-B_i(0,t)],\qquad \tau_A(t)=\mathbb E_P[A_i(1,t)-A_i(0,t)].\]
Specify the information benchmark \(b^*(I)=\mathbb E[R\mid I]\) and examine whether the treatment changes forecast errors relative to it. A motivation interpretation requires ruling out treatment-induced information, altered forecast incentives, or rational inferences about the experiment. To determine whether distorted beliefs cause the action effect, additionally randomize a credible belief-correction intervention or justify an explicit mediation model; correlation between \(B_i\) and \(A_i\) is insufficient. Estimate persistence over preregistered horizons and heterogeneity across initially winning and losing positions. The broader question is whether these identified mechanisms recur with economically important magnitudes in consequential saving and investment decisions, under comparable information and ownership conditions. It should not be treated as a theorem that all optimistic beliefs are motivated or that the stock-market findings already identify population prevalence.

\emph{Definitions and maintained assumptions (editorial unless a source definition is named).} A full objective probability law for \((R,I)\) is needed to define \(b^*(I)=\mathbb E[R\mid I]\); prices and publicly observed information alone do not provide an objectively known expected return in an actual stock market. In a controlled experiment choose that law explicitly, or call a field benchmark a maintained forecast model. Treatment \(Z\) specifies ownership, endowment or self-image incentives and what it changes. Forecast elicitation has a fixed scoring rule and known stakes, since treatment may otherwise change forecast incentives. Potential beliefs and actions are measured at specified horizons before any treatment-dependent stopping or sale; selection into continued ownership must be handled explicitly. The belief-error effect is \(\mathbb E[(B(1,t)-b^*(I))-(B(0,t)-b^*(I))]\) under unchanged information. A second intervention targeting beliefs identifies an action response only if its direct utility/information effects are excluded. Winning/losing subgroups should be defined by preassigned exogenous price paths, not treatment-induced portfolios.
\subsection*{Variants and assumption changes}
\begin{enumerate}
\item \textbf{Known-return ownership (editorial).} Randomize stock endowment under an announced finite return distribution and identical forecast payoffs. Estimate ownership effects on forecast error and sale decisions.
\item \textbf{Field persistence (editorial).} Use comparable real equities and a preregistered forecasting benchmark over repeated horizons. Estimate treatment effects and correction responses with explicit attrition and information assumptions; a benchmark-dependent optimism measure is not proof of irrationality.
\end{enumerate}
\subsection*{References and source audit}
\begin{enumerate}
\item Publisher-indexed and author abstract verify article/results; full publisher paper inaccessible. No exact broad open-question locator. \url{https://doi.org/10.1016/j.jbankfin.2025.107510}
\item Publisher-indexed and author abstract verify article/results; full publisher paper inaccessible. No exact broad open-question locator. \url{https://sites.google.com/site/carloscuevaherrero/home/research}
\end{enumerate}
\begin{enumerate}\raggedright
\item\hypertarget{definitions:source:30007587:1}{} Carlos Cueva; Iñigo Iturbe-Ormaetxe. 2025. \emph{Motivated beliefs about stock returns}. Abstract, corroborated by author research page https://sites.google.com/site/carloscuevaherrero/home/research. Full paper was not accessible in this audit..
\url{https://doi.org/10.1016/j.jbankfin.2025.107510}.
DOI: \url{https://doi.org/10.1016/j.jbankfin.2025.107510}.
\item\hypertarget{definitions:source:30007587:2}{} Carlos Cueva. 2025. \emph{Author research page: Motivated beliefs about stock returns}. Published, entry Motivated beliefs about stock returns.
\url{https://sites.google.com/site/carloscuevaherrero/home/research}.
\end{enumerate}

\subsection*{Common conventions: Source mathematical conventions}

These conventions apply to each entry unless explicitly replaced.

\textbf{Models and identification.} A primitive model \(m\) specifies all probability laws, feasible choices, information, equilibrium and measurement rules used by its target. The admissible class \(\mathcal M\) is fixed independently of outcomes. If \(P_O(m)\) is its observable probability law and \(\tau(m)\) the target, its identified set at observable law \(P\) is
\[
 \mathcal I_\tau(P)=\{\tau(m):m\in\mathcal M,\ P_O(m)=P\}.
\]
Point identification means that this set is a singleton. An empty set signals incompatibility of data and assumptions. Coordinate bounds need not describe the joint set. Bounds are sharp only if every retained target is attainable or arbitrarily approachable by compatible models. Finite bounds require explicit restrictions on unobserved quantities; unrestricted confounding, hidden wealth, extrapolation or arbitrary equilibrium selection can yield uninformative sets.

\textbf{Causal interventions.} A potential outcome \(Y_i(p)\) is the outcome under a complete policy regime \(p\), including timing, financing, information and market spillovers. The target population and units are fixed before treatment. With interference, use the complete assignment vector or a justified exposure mapping; individual no-interference notation alone cannot identify an economy-wide intervention. A difference of mean potential outcomes does not require the cross-world joint distribution. Conditioning on post-treatment migration, survival, enrollment or employment changes the estimand and needs separate justification.

\textbf{Statistical statements.} An estimator, confidence set, forecast or test specifies a sampling law, model class and loss/coverage criterion. Uniform \((1-\alpha)\)-coverage means \(\inf_{m\in\mathcal M}\Pr_m\{\tau(m)\in C_n\}\geq1-\alpha\), or an explicitly asymptotic version. The whole parameter space trivially has coverage, so informativeness requires a stated width, risk or power objective. A forecasting improvement is relative to a named benchmark, information set and evaluation population.

\textbf{Equilibrium and optimization.} An equilibrium comprises feasible plans, optimality, beliefs where needed, budget constraints and all market/resource clearing. The primitives or a stated benchmark define each correspondence \(\mathcal E(p;m)\). Nonemptiness is assumed or proved, never inferred just from compact policy sets. A selected-equilibrium target uses a declared selection rule; a robust target quantifies over all permitted equilibria. Use suprema or infima unless attainment is established. An \(\varepsilon\)-optimal policy is feasible and within \(\varepsilon\) of the finite optimum value. Report an empty feasible set separately.

\textbf{Welfare and accounting.} Normative utility, interpersonal normalization, social weights, numeraire, discounting and the use of public receipts are primitives. Behavioral utility need not coincide with the welfare criterion. Household welfare includes ownership income and rents once; adding profits again to compensating variation generally double-counts them. A monetary equivalent is defined by an explicit transfer and price convention, with monotonicity and range conditions. Average effects, mechanism contributions, transfers and welfare losses are different targets. Infinite sums require convergence and dynamic feasible plans require terminal or no-Ponzi conditions.

\textbf{Source versions.} A working-paper date, revision date and journal publication year are distinguished. A DOI for a journal version is not automatically the DOI of an earlier manuscript. Locators refer to the stated version and printed pages where specified. A metadata-only check does not verify a claim about a full-text conclusion. Every entry's source subsection narrows the provenance claim accordingly.
% END ECONOMICS STATEMENT
\end{document}
